% Die Titelseite der Arbeit

\begin{titlepage}

\begin{center} % zentrieren

  % Logo der Universit{\"a}t Mannheim
  \begin{figure}[ht]
    \centering
    \includegraphics[width=.6\textwidth]{grafiken/unilogo}
  \end{figure}
  
  % Vertikaler Zwischenraum
  \bigskip
  \vfill 
  \begin{framed}
  % Titel der Arbeit und Typ der Arbeit, umrandet
    \begin{center}
     \textsc{{\Large Seizing Continuous Integration Process to Improve Software Reuse}}
                                % Letztes \\ ist wichtig, beginnt eine neue Zeile f{\"u}r die Art der Arbeit
  
      \bigskip
  
                                % Art der Arbeit, ggf. auszutauschen gegen Seminar- oder Doktorarbeit
      \textbf{Master Thesis}
    \end{center}
    \end{framed}
    \vfill
    \vfill
  
  % Daten des Erstellers, Einreichungsdatum
  % in einer Tabelle ausgerichtet
  \begin{tabular*}{0.62\textwidth}{r@{\extracolsep{\fill}}l}
   submitted: &\ February 2017\\\\
    by: &\ Piero Antonio Divasto Martínez\\
    &\ born March 10th, 1986\\
    &\ in Iquique\\
    &\ Chile\\
    \\
    Student ID Number: &\ 1472793\\
  \end{tabular*}
  \vfill
  \vfill
  
  % Unten: Kontaktdaten des Lehrstuhls f{\"u}r Wirtschaftsinformatik 1
  
  \rule{\textwidth}{.4pt}\\ % vertikale Linie
  University of Mannheim\\
  Chair of Software Engineering\\
  B6, 26, Gebäudeteil C\\
  D -- 68159 Mannheim\\
  Phone: +49 621-181-1691, Fax +49 621-181-1692\\
  Internet: \url{http://swt.informatik.uni-mannheim.de/}
\end{center}

\end{titlepage} % Ende des Titelblatts

%%% Local Variables: 
%%% mode: latex
%%% TeX-master: "~/Documents/DA-Vorlage/beispiel/da-beispiel"
%%% End: 
